\documentclass[12pt,a4paper]{article}

% ---------------------------------------------------------------
% Pacotes
% ---------------------------------------------------------------
\usepackage[utf8]{inputenc}
\usepackage[T1]{fontenc}
\usepackage[brazil]{babel}
\usepackage{lmodern}
\usepackage{microtype}
\usepackage[top=3cm,bottom=2cm,left=3cm,right=2cm]{geometry}
\usepackage{setspace}
\usepackage{indentfirst}
\usepackage{enumitem}
\usepackage{booktabs}
\usepackage{tabularx}
\usepackage{array}
\usepackage{xcolor}
\usepackage{amsmath}
\usepackage{titlesec}
\usepackage{fancyhdr}
\usepackage[hidelinks]{hyperref}
\usepackage{xurl}
\usepackage{etoolbox}
\patchcmd{\thebibliography}{\section*{\refname}}{\subsection*{\refname}}{}{}
\usepackage{float}

\onehalfspacing
\setlength{\parindent}{1.25cm}
\setlength{\parskip}{0.3em}

\definecolor{caixa}{gray}{0.95}

\titleformat{\section}{\large\bfseries}{\thesection}{0.6em}{\MakeUppercase}
\titleformat{\subsection}{\normalsize\bfseries}{\thesubsection}{0.6em}{}

\pagestyle{fancy}
\fancyhf{}
\fancyhead[L]{\small Engenharia de Software -- Trabalho Final}
\fancyhead[R]{\small\thepage}
\renewcommand{\headrulewidth}{0.4pt}

% Caixa para o texto de cada história de usuário
\newcommand{\historia}[1]{%
  \par\noindent\colorbox{caixa}{%
    \parbox{\dimexpr\linewidth-2\fboxsep\relax}{\itshape #1}}\par\medskip}

% Lista de critérios numerada como 1.1, 1.2, ...
\newlist{criterios}{enumerate}{1}
\setlist[criterios]{leftmargin=1.4cm, labelsep=0.4cm, itemsep=0.4em}

\newcommand{\repositorio}{https://github.com/GustavoFerreiraFonseca/ProjetoFinal_Engenharia_Software}

% ---------------------------------------------------------------
\begin{document}

% ---------------------------------------------------------------
% Capa
% ---------------------------------------------------------------
\begin{titlepage}
  \centering
  {\large UNIVERSIDADE TECNOLÓGICA FEDERAL DO PARANÁ\par}
  \vspace{0.3cm}
  {\large Engenharia de Software\par}

  \vspace{2.5cm}
  {\normalsize Gustavo Ferreira da Fonseca\par}
  {\normalsize Gustavo Henrique Gonçalves\par}
  {\normalsize João Francisco Barcala Paulo\par}

  \vfill
  {\LARGE\bfseries Analisador de Desempenho e Rotinas de Treino para Corrida\par}
  \vspace{0.8cm}
  {\large Documentação do Trabalho Final\par}
  \vfill

  {\normalsize 2026\par}
\end{titlepage}

\tableofcontents
\newpage

% ---------------------------------------------------------------
\section{Definição do trio}
% ---------------------------------------------------------------

\begin{table}[H]
  \centering
  \begin{tabular}{lc}
    \toprule
    \textbf{Nome completo} & \textbf{RA} \\
    \midrule
    Gustavo Ferreira da Fonseca  & 2669510 \\
    Gustavo Henrique Gonçalves   & 2669528 \\
    João Francisco Barcala Paulo & \\
    \bottomrule
  \end{tabular}
\end{table}

\noindent\textbf{Repositório:} \url{\repositorio}

\subsection{Acompanhamento das sprints}

\begin{table}[H]
  \centering
  \begin{tabularx}{\linewidth}{l c X}
    \toprule
    \textbf{Sprint} & \textbf{Review} & \textbf{Entrega} \\
    \midrule
    Sprint 1 & 30/10 & Definição do trio, descrição do sistema e pesquisa (Seções 1 a 3) \\
    Sprint 2 & 06/11 & Histórias de usuário, critérios de aceitação e testes (Seções 4 e 5) \\
    Sprint 3 & 13/11 & Diagrama da arquitetura e desenvolvimento das histórias (Seção 6) \\
    Sprint 4 & 27/11 & Refatoração e padrões de projeto (Seção 7) \\
    \bottomrule
  \end{tabularx}
\end{table}

% ---------------------------------------------------------------
\section{Descrição do sistema}
% ---------------------------------------------------------------

\noindent\textbf{Tema:} Analisador de Desempenho e Rotinas de Treino para Corrida.

Esse será um sistema em que corredores poderão registrar seus treinos, além de
acompanhar sua evolução ao longo das semanas de uso. Primeiramente, receberemos um
arquivo CSV com todo o histórico de treino (data, tempo, distância, frequência
cardíaca). Pegamos esses dados, processaremos para obter pace médio, volume semanal
acumulado e recordes pessoais. Além disso, iremos realizar análises para verificar se
não houve um aumento abrupto de volume que pode gerar lesões. Por fim, gerará um
relatório em txt consolidando sua evolução semanal, destaques de desempenho e alertas
de risco de lesões.

\begin{table}[H]
  \centering
  \caption{Colunas do arquivo CSV de entrada.}
  \label{tab:csv}
  \begin{tabularx}{\linewidth}{l X}
    \toprule
    \textbf{Coluna} & \textbf{Descrição} \\
    \midrule
    \texttt{data}          & Data do treino, no formato DD/MM/AAAA \\
    \texttt{distancia\_km} & Distância percorrida, em quilômetros \\
    \texttt{tempo\_min}    & Tempo total do treino, em minutos \\
    \texttt{fc\_media}     & Frequência cardíaca média do treino, em bpm \\
    \bottomrule
  \end{tabularx}
\end{table}

% ---------------------------------------------------------------
\section{Pesquisa e justificativa}
% ---------------------------------------------------------------

\noindent\textbf{Fontes de referência:}

\begin{itemize}[leftmargin=0.6cm, itemsep=0.5em]
  \item[] CAI, Jie; LI, Gang. Exercise or lie down? The impact of fitness app use on
        users' wellbeing. Frontiers in Public Health, v. 11, p. 1281323, 2024. DOI:
        10.3389/fpubh.2023.1281323.
  \item[] VARELLA, Drauzio. Correr: O exercício, a cidade e a busca da superação. 1. ed.
        São Paulo: Companhia das Letras, 2015.
  \item[] LABORO. Os benefícios da prática de atividade física orientada via aplicativo.
        Blog Grupo Laboro, 2021. Disponível em:
        \url{https://laboro.edu.br/blog/os-beneficios-da-pratica-de-atividade-fisica-orientada-via-aplicativo/}.
        Acesso em: 30 out. 2026.
  \item[] VARELLA, Drauzio. Atividade física e metabolismo. Portal Drauzio Varella, UOL.
        Disponível em:
        \url{https://drauziovarella.uol.com.br/drauzio/atividade-fisica-e-metabolismo-artigo/}.
        Acesso em: 30 out. 2026.
\end{itemize}

As referências acima sustentam a importância da atividade física e, mais
especificamente, os artigos ``Exercise or lie down? The impact of fitness app use on
users' wellbeing. Frontiers in Public Health'' e ``Os benefícios da prática de
atividade física orientada via aplicativo. Blog Grupo Laboro, 2021'' sustentam que o
monitoramento da evolução por meio de aplicativos podem auxiliar os praticantes a
começarem e a manterem a disciplina nos exercícios físicos.

% ---------------------------------------------------------------
\section{Histórias de usuário e critérios de aceitação}
% ---------------------------------------------------------------

\subsection{História de usuário 1}

\historia{1 - Como usuário eu quero carregar um arquivo csv com meu histórico de
corrida e obter o cálculo do meu pace médio por treino para acompanhar minha
evolução nos meus treinos.}

\noindent\textbf{Critérios de aceitação:}
\begin{criterios}[label=\textbf{1.\arabic*:}]
  \item O sistema deve ler o arquivo csv e ler a distância e o tempo total de cada
        corrida. Em quilômetros e minutos respectivamente.
  \item Para cada corrida registrada, o sistema deve calcular corretamente o pace em
        minutos por quilômetro, dividindo o tempo decorrido pela distância, se
        atentando a conversão da parte decimal para segundos.
  \item Se o arquivo csv contiver distância igual ou inferior a 0 ou tempo igual ou
        inferior a 0, o sistema deve ignorar essa linha.
  \item O pace resultante deve ser formatado seguindo o seguinte exemplo: 5:30 min/km.
\end{criterios}

\subsection{História de usuário 2}

\historia{2 - Como usuário eu quero agrupar os meus treinos semanalmente, e a partir
disso calcular o volume de treino e meus recordes.}

\noindent\textbf{Critérios de aceitação:}
\begin{criterios}[label=\textbf{2.\arabic*:}]
  \item O sistema, ao receber os dados dos treinos, deve organizá-los por semana e
        ordená-los por data caso não esteja ordenado.
  \item Ao analisar os treinos, o sistema irá calcular a somatória das distâncias
        percorridas na semana para determinar o volume semanal.
  \item O sistema deverá identificar e atualizar os recordes pessoais do usuário, que
        possui 4 campos de recorde pessoal:
        \begin{itemize}[itemsep=0pt]
          \item Maior distância (km).
          \item Maior volume semanal acumulado (km).
          \item Maior número de treinos realizados em uma semana.
          \item Menor pace médio registrado em um treino (min/km).
        \end{itemize}
  \item Cada agrupamento semanal deve ser identificado pelo primeiro dia da semana
        (Domingo), seguindo com o padrão DD/MM/AAAA. Como exemplo: ``semana
        01/11/2026:''.
\end{criterios}

\subsection{História de usuário 3 -- Comparativo semanal e disparo de alertas de lesão}

\historia{3 - Como corredor, eu quero receber alertas quando meu volume semanal, meu
ritmo ou a sequência dos meus treinos mudarem de forma brusca, para reduzir o risco de
ter uma lesão.}

A história usa três alertas. O alerta de volume usa o ACWR (\textit{Acute:Chronic
Workload Ratio}) de Gabbett \cite{gabbett2016}, que compara o volume recente do corredor
(agudo, a semana atual) com o volume com que ele já está acostumado (crônico, a média
das últimas 4 semanas):
\[
  \text{ACWR} = \frac{\text{volume agudo}}{\text{volume crônico}}
\]

Por exemplo, um corredor que fez 10, 10 e 10 km nas três semanas anteriores e 16 km na
semana atual tem volume crônico de $(10 + 10 + 10 + 16) / 4 = 11{,}5$ km e ACWR de
$16 / 11{,}5 = 1{,}39$, o que gera um alerta de atenção (Tabela~\ref{tab:acwr}).

O alerta de ritmo vem de Nielsen \textit{et al.} \cite{nielsen2013}, que ligam mudanças
bruscas de pace a lesões no tendão de Aquiles, na panturrilha e na fáscia plantar. O
alerta de recuperação avisa quando o corredor treina vários dias seguidos sem descanso.

\noindent\textbf{Critérios de aceitação:}
\begin{criterios}[label=\textbf{3.\arabic*:}]
  \item Para cada semana (agrupamento da história 2), o sistema deve calcular o volume
        agudo (soma das distâncias da semana), o volume crônico (média do volume das
        últimas 4 semanas, incluindo a atual; semanas sem treino contam como 0 km) e o
        ACWR, arredondado para 2 casas decimais.
  \item O sistema deve classificar o ACWR seguindo as faixas da
        Tabela~\ref{tab:acwr}. As mensagens de alerta devem seguir os exemplos:
        \begin{itemize}[itemsep=2pt, before=\raggedright\small]
          \item \texttt{ATENÇÃO (volume): 16.0 km nesta semana contra média de 11.5 km nas últimas 4 semanas (ACWR 1.39).}
          \item \texttt{ALERTA (volume): 30.0 km nesta semana contra média de 15.0 km nas últimas 4 semanas (ACWR 2.00). Zona de perigo para lesões.}
        \end{itemize}
  \item O sistema deve calcular o pace de cada treino (tempo / distância), o pace médio
        da semana (média dos paces dos treinos da semana) e o pace crônico (média dos
        paces dos treinos das últimas 4 semanas, incluindo a atual). Se o pace da
        semana for 15\% ou mais rápido que o crônico, pela fórmula
        \[
          \text{diferença} = \frac{\text{pace crônico} - \text{pace da semana}}
                                  {\text{pace crônico}} \times 100
        \]
        arredondada para 1 casa decimal, o sistema deve gerar o alerta:
        \begin{itemize}[itemsep=2pt, before=\raggedright\small]
          \item \texttt{ALERTA (ritmo): pace médio da semana de 4:48 min/km, 15.8\% mais rápido que a média das últimas 4 semanas (5:42 min/km).}
        \end{itemize}
  \item Para cada sequência de 3 ou mais dias seguidos com treino, o sistema deve gerar
        um único alerta de recuperação com o primeiro e o último dia da sequência. Dois
        treinos no mesmo dia contam como um dia só.
        \begin{itemize}[itemsep=2pt, before=\raggedright\small]
          \item \texttt{ALERTA (recuperação): treinos em 3 dias seguidos, de 02/11/2026 a 04/11/2026. Agende um dia de descanso.}
        \end{itemize}
  \item Se a semana não tiver 4 semanas de histórico (contando a partir da primeira
        semana com treino), o sistema não calcula o ACWR nem a diferença de pace,
        classifica a semana como ``Sem base de comparação'' e não gera alertas de volume
        nem de ritmo.
\end{criterios}

\begin{table}[H]
  \centering
  \caption{Faixas de classificação do ACWR (critério 3.2), com base em
           Gabbett \cite{gabbett2016}.}
  \label{tab:acwr}
  \begin{tabular}{lll}
    \toprule
    \textbf{ACWR} & \textbf{Classificação} & \textbf{Alerta} \\
    \midrule
    Abaixo de 0,8              & Abaixo da zona segura   & Não \\
    De 0,8 a 1,3               & Zona segura             & Não \\
    Acima de 1,3 até 1,5       & Atenção                 & Sim \\
    Acima de 1,5               & Zona de perigo          & Sim \\
    Sem 4 semanas de histórico & Sem base de comparação  & Não \\
    \bottomrule
  \end{tabular}
\end{table}

% ---------------------------------------------------------------
\section{Relação entre testes e critérios de aceitação}
% ---------------------------------------------------------------

Os testes automatizados estão na pasta \texttt{testes/}, com um arquivo para cada
história de usuário, e usam o módulo \texttt{unittest} do Python. O nome de cada teste
começa pelo número do critério que ele verifica. Para rodar todos os testes a partir da
raiz do repositório:

\begin{verbatim}
python -m unittest discover -s testes -v
\end{verbatim}

\begin{table}[H]
  \centering
  \caption{Relação entre os testes automatizados e os critérios de aceitação.}
  \label{tab:testes}
  \small
  \begin{tabularx}{\linewidth}{l >{\ttfamily}X c}
    \toprule
    \textbf{Arquivo} & \textbf{\normalfont Teste} & \textbf{Critério(s)} \\
    \midrule
    \texttt{testes\_historia1.py} & test\_1\_1\_leitura\_dados\_validos             & 1.1 \\
                                  & test\_1\_2\_e\_1\_4\_calculo\_e\_formatacao\_pace   & 1.2 e 1.4 \\
                                  & test\_1\_3\_ignorar\_linhas\_invalidas          & 1.3 \\
    \midrule
    \texttt{testes\_historia2.py} & test\_2\_1\_e\_2\_4\_agrupamento\_e\_identificacao & 2.1 e 2.4 \\
                                  & test\_2\_2\_calculo\_volumeSemanal              & 2.2 \\
                                  & test\_2\_3\_recordePessoal                      & 2.3 \\
    \midrule
    \texttt{testes\_historia3.py} & test\_3\_1\_volume\_semanal\_e\_acwr             & 3.1 \\
                                  & test\_3\_2\_classificacao\_e\_alertas\_de\_volume                & 3.2 \\
                                  & test\_3\_3\_mudanca\_brusca\_de\_pace                & 3.3 \\
                                  & test\_3\_4\_dias\_seguidos\_sem\_descanso     & 3.4 \\
                                  & test\_3\_5\_semanas\_sem\_historico\_suficiente               & 3.5 \\
    \bottomrule
  \end{tabularx}
\end{table}

% ---------------------------------------------------------------
\section{Arquitetura do sistema}
% ---------------------------------------------------------------

\textit{Seção a ser desenvolvida na Sprint 3: diagrama da arquitetura, com os
componentes, suas responsabilidades e a comunicação entre eles.}

% ---------------------------------------------------------------
\section{Padrões de projeto}
% ---------------------------------------------------------------

\textit{Seção a ser desenvolvida na Sprint 4: padrões de projeto aplicados, com a
justificativa de cada um e um diagrama por padrão com as classes e módulos reais do
código.}

% ---------------------------------------------------------------
% Referências (Sprint 1 e história de usuário 3)
% ---------------------------------------------------------------
\newpage
\section*{Referências}
\addcontentsline{toc}{section}{Referências}

\renewcommand{\refname}{Pesquisa (Sprint 1)}
\begin{thebibliography}{9}

\bibitem{cai2024}
CAI, Jie; LI, Gang. Exercise or lie down? The impact of fitness app use on users'
wellbeing. \textbf{Frontiers in Public Health}, v. 11, p. 1281323, 2024.
DOI: 10.3389/fpubh.2023.1281323.

\bibitem{varella2015}
VARELLA, Drauzio. \textbf{Correr}: o exercício, a cidade e a busca da superação.
1. ed. São Paulo: Companhia das Letras, 2015.

\bibitem{laboro2021}
LABORO. Os benefícios da prática de atividade física orientada via aplicativo.
\textbf{Blog Grupo Laboro}, 2021. Disponível em:
\url{https://laboro.edu.br/blog/os-beneficios-da-pratica-de-atividade-fisica-orientada-via-aplicativo/}.
Acesso em: 30 out. 2026.

\bibitem{varella_metabolismo}
VARELLA, Drauzio. Atividade física e metabolismo. \textbf{Portal Drauzio Varella},
UOL. Disponível em:
\url{https://drauziovarella.uol.com.br/drauzio/atividade-fisica-e-metabolismo-artigo/}.
Acesso em: 30 out. 2026.

\end{thebibliography}

\renewcommand{\refname}{História de usuário 3}
\begin{thebibliography}{9}
\setcounter{enumiv}{4}

\bibitem{gabbett2016}
GABBETT, Tim J. The training-injury prevention paradox: should athletes be training
smarter and harder? \textbf{British Journal of Sports Medicine}, v. 50, n. 5,
p. 273-280, 2016. DOI: 10.1136/bjsports-2015-095788.

\bibitem{nielsen2013}
NIELSEN, Rasmus O. \textit{et al.} Classifying running-related injuries based upon
etiology, with emphasis on volume and pace. \textbf{International Journal of Sports
Physical Therapy}, v. 8, n. 2, p. 172-179, 2013.

\end{thebibliography}

\end{document}
